%=====================================================================
%  Part 4 opener --- Measurement, Data and Analysis
%=====================================================================
\structuralfigures{P4.}

A design produces numbers or text; neither speaks for itself. This part
covers the two halves of turning observation into a defensible conclusion:
making a construct observable without distorting it, and reasoning from the
result without claiming more than it supports. The statistics here are chosen
for what graduate computing research actually uses, and every technique is
introduced with the mistake it is most often used to commit.

\begin{conceptbox}[What Part~4 Covers]
\begin{tight}
  \item \textbf{Chapter 17 --- Measurement and Instrument Validity.} Scales, construct validity, reliability, and why Cronbach's alpha is not the reassurance it appears to be.
  \item \textbf{Chapter 18 --- Sampling and Study Power.} Probability and purposive sampling, a-priori power, and the arithmetic that decides feasibility.
  \item \textbf{Chapter 19 --- Statistical Inference Done Properly.} What a p-value is not, effect sizes with intervals, assumptions, and multiple comparisons.
  \item \textbf{Chapter 20 --- Regression and Modelling.} Linear, logistic and mixed-effects models; inference versus prediction; interpreting coefficients honestly.
  \item \textbf{Chapter 21 --- Structural Equation Modelling.} Latent constructs, CB-SEM and PLS-SEM, discriminant validity and common method bias.
  \item \textbf{Chapter 22 --- Evaluating Machine Learning Research.} Leakage, validation that matches deployment, comparing classifiers, and documenting models and data.
  \item \textbf{Chapter 23 --- Qualitative Data Analysis in Practice.} Coding with open-source tools, the audit trail, and trustworthiness criteria.
  \item \textbf{Chapter 24 --- Threats to Validity.} The four-way taxonomy applied to computing, and how to write the section examiners look for.
\end{tight}
\end{conceptbox}
